\chapter{Variables de entorno}
\label{sec:envvars}

\section{Convención}

Todas las variables del producto tienen prefijo
\code{GRAPHHUNTER_}. La precedencia es:

\begin{enumerate}
    \item Variable del proceso (seteada por el OS o por el shell).
    \item \filepath{.env} en el working dir del Tauri daemon.
    \item Default hardcodeado en código.
\end{enumerate}

Un proceso arrancado con una variable explícita ignora
\filepath{.env}. Esto es deliberado: en producción
\filepath{.env} no existe y las variables vienen del supervisor.

\section{Daemon (Tauri app)}

\begin{table}[H]
\centering
\small
\begin{tabular}{@{}p{5cm}lp{5cm}@{}}
\toprule
Variable & Default & Propósito \\
\midrule
\code{GRAPHHUNTER_API_TOKEN} & autogenerado & Bearer token para la
API HTTP. Impreso al arrancar; copiarlo al cliente MCP. \\
\code{GRAPHHUNTER_API_PORT} & \code{37891} & Puerto TCP. \\
\code{GRAPHHUNTER_API_HOST} & \code{127.0.0.1} & Host de bind. Usar
\code{0.0.0.0} sólo en entornos confiables. \\
\code{GRAPHHUNTER_CATALOG_DIR} & \code{\$HOME/.graphhunter/catalog} &
Directorio con YAMLs de catálogo custom. \\
\code{GRAPHHUNTER_LOCAL_MODEL_PATH} & (ninguno) & Ruta a modelo GGUF
para backend local-llm. \\
\code{GRAPHHUNTER_IOC_PUBLISH_ENABLED} & \code{0} & Habilita
\code{POST /publish_iocs_to_sentinel}. \\
\code{GRAPHHUNTER_TICKET_PUBLISH_ENABLED} & \code{0} & Habilita
\code{POST /file_ticket}. \\
\bottomrule
\end{tabular}
\caption{Variables consumidas por el daemon.}
\end{table}

\section{Cliente MCP}

\begin{table}[H]
\centering
\small
\begin{tabular}{@{}p{5.5cm}p{3cm}p{4cm}@{}}
\toprule
Variable & Default & Propósito \\
\midrule
\code{GRAPHHUNTER_API_URL} & \code{http://127.0.0.1:37891} & URL del
daemon. \\
\code{GRAPHHUNTER_API_TOKEN} & (obligatorio) & Mismo token que el
daemon imprime. \\
\code{GRAPHHUNTER_MCP_DEBUG} & \code{0} & Log verboso. \\
\code{GRAPHHUNTER_MCP_ALLOW_NO_TOKEN} & \code{0} & \emph{Sólo dev}:
permite arrancar sin token. \\
\code{GRAPHHUNTER_CONTRACT_TEST} & \code{0} & Habilita suite de
tests contractuales. \\
\bottomrule
\end{tabular}
\caption{Variables consumidas por el servidor MCP.}
\end{table}

\section{Integraciones externas}

Consumidas por \code{POST /file_ticket}:

\begin{itemize}
    \item \textbf{Jira}: \code{JIRA_BASE_URL}, \code{JIRA_EMAIL},
    \code{JIRA_API_TOKEN}, \code{JIRA_PROJECT_KEY},
    \code{JIRA_ISSUE_TYPE} (default \code{Task}).
    \item \textbf{ServiceNow}: \code{SERVICENOW_INSTANCE},
    \code{SERVICENOW_USERNAME}, \code{SERVICENOW_PASSWORD},
    \code{SERVICENOW_TABLE} (default \code{incident}).
\end{itemize}

Consumidas por \code{POST /publish_iocs_to_sentinel}:
credenciales de Azure (\code{AZURE_TENANT_ID},
\code{AZURE_CLIENT_ID}, \code{AZURE_CLIENT_SECRET}) y el workspace
ID de Sentinel (\code{SENTINEL_WORKSPACE_ID}).

\section{Feature flags}

Estas dos flags son las que cambian superficie expuesta:

\begin{itemize}
    \item \code{GRAPHHUNTER_IOC_PUBLISH_ENABLED=1}: sin esto,
    \code{/publish_iocs_to_sentinel} responde 403.
    \item \code{GRAPHHUNTER_TICKET_PUBLISH_ENABLED=1}: análogo para
    \code{/file_ticket}.
\end{itemize}

Cada tool MCP con efectos externos también pide \code{confirm:true}
por llamada, o \code{dry_run:true} para previsualizar, sobre los
env vars anteriores. Doble barrera.

\section{Flags de \emph{performance}}

Las siguientes flags activan rutas de \emph{matcher} optimizadas
incorporadas en el \emph{M-track} de abril 2026. Todas son
\emph{opt-in} y compatibles con la API existente: cuando están
desactivadas el comportamiento es el del \emph{path} legacy.

\begin{table}[H]
\centering
\small
\begin{tabular}{@{}p{6cm}p{2cm}p{6cm}@{}}
\toprule
Variable & Default & Propósito \\
\midrule
\code{GRAPHHUNTER_K4_LFTJ} & \code{0} & Habilita el detector \emph{name marker} S3 que
rutea hipótesis con \emph{shape K\_4} hacia el \emph{enumerador} LFTJ
\code{search\_temporal\_pattern\_k4\_lftj}. \\
\code{GRAPHHUNTER_LFTJ_PLANNER} & \code{0} & Habilita el detector
\emph{estructural} M1 (6 pasos homogéneos sin predicados) que rutea a
LFTJ-K4 sin requerir \emph{name marker}. Ortogonal a
\code{K4\_LFTJ}; ambos pueden coexistir. \\
\code{GRAPHHUNTER_YANNAKAKIS} & \code{0} & L1: habilita el
\emph{matcher} de cadenas \emph{output-sensitive} (semi-join
reducer de dos pasadas + \emph{path enumerator}). Aplica a toda
hipótesis $\alpha$-acíclica --- las cadenas del DSL actual lo son.
Cae al DFS legacy cuando el \emph{planner} detecta que la
hipótesis tiene un ciclo (vía \emph{shared bindings}) o cuando el
\emph{flag} está apagado. Default-on en los binarios Tauri/CLI;
override con \code{GRAPHHUNTER_YANNAKAKIS=0}. \\
\code{GRAPHHUNTER_AMAC} & \code{0} & Activa el \emph{DFS} \emph{interleaved} estilo AMAC
con $K=8$ contextos en vuelo por \emph{worker} de Rayon, pensado para
\emph{hide} latencia de DRAM en patrones $\geq 3$ saltos. \\
\code{GRAPHHUNTER_SEMIJOIN} & \code{0} & Activa la pasada
\emph{semi-join} bidireccional D2-lite que poda \emph{candidate sets}
antes de entrar al DFS. \\
\code{GRAPHHUNTER_HUGEPAGES} & \code{0} & Sólo Linux: invoca
\code{madvise(MADV\_HUGEPAGE)} sobre regiones \code{Mmap} para
reducir \emph{TLB stalls} en grafos $> L2$. \\
\code{GRAPHHUNTER_MEM_WARN} & \code{0.75} & Heavy-logs:
\emph{ratio} \(\in (0, 1]\) sobre el cual la \emph{pildora} \emph{Mem:}
del \emph{toolbar} pasa a \emph{amber}. Usado por la UI para señalar
presión de memoria moderada. \\
\code{GRAPHHUNTER_MEM_CRITICAL} & \code{0.90} & Heavy-logs:
análogo a \code{MEM\_WARN} pero para el estado \emph{red}. Por
encima de este \emph{ratio} la \emph{pildora} aparece en rojo y la
sesión está cerca del límite de la host RAM. \\
\code{GRAPHHUNTER_INDEX_HEAP_BUDGET} & \code{1073741824} & Heavy-logs:
\emph{budget} en \emph{bytes} sobre el cual los \emph{auxiliary
indexes} (NLF, k-hop reachability) se \emph{spillan} a un \emph{temp
file} \emph{mmap-backed} en lugar de quedar en \emph{heap}. Default
$= 1$~GiB. Sobre un EVTX de 28~GB cada \emph{array} temporal de
\code{KHopReach} pesa $\sim$2.16~GB y por lo tanto se \emph{spilla}
automáticamente, manteniendo el RSS acotado por el \emph{working set}
en lugar de por el tamaño total del \emph{log}. Setear a \code{0}
fuerza \emph{spill} siempre; valores altos fuerzan \emph{heap}. La
selección sucede \emph{per-index} al construirse, sin \emph{user
config} requerida. \\
\bottomrule
\end{tabular}
\caption{Variables consumidas por la capa de \emph{matcher}. Sin
ellas el comportamiento es el legacy DFS.}
\end{table}

\textbf{Capability de simetría \emph{nauty}.} El módulo
\filepath{core/graph-engine/src/nauty\_symmetry.rs} expone una
constante \code{HAS\_NAUTY: bool} que vale \code{true} sólo cuando
el \emph{crate} se compiló con \code{--features symmetry-nauty}
sobre un \emph{target} \code{cfg(unix)}. En Windows o sin la
\emph{feature} la \emph{capability} es \code{false} y el módulo
retorna \emph{constraints} vacíos como \emph{soundness-preserving
no-op}. No es una variable de entorno --- es una decisión de
\emph{build}, descrita en \filepath{core/graph-engine/Cargo.toml}.

\section{Fuente canónica}

El chequeo de env vars ocurre en:

\begin{itemize}
    \item \filepath{apps/tauri/src-tauri/src/lib.rs}: token, port.
    \item \filepath{apps/tauri/src-tauri/src/http/mod.rs}: host.
    \item \filepath{core/graph-engine/src/catalog/loader.rs}:
    catalog dir.
    \item \filepath{platform/local-llm/src/lib.rs}: model path.
    \item \filepath{graph-hunter-mcp/src/index.ts}: cliente MCP.
\end{itemize}

Agregar una variable nueva obliga a actualizar este capítulo y
\filepath{.env.example}.
